% Generated draft: requires mathematical and source review.
\documentclass[11pt]{article}
\usepackage[margin=25mm]{geometry}
\usepackage{fontspec}
\setmainfont{Latin Modern Roman}
\usepackage{amsmath,amssymb,mathtools,unicode-math}
\setmathfont{Latin Modern Math}
\usepackage{enumitem,needspace,xurl,hyperref}
\setlength{\parindent}{0pt}
\setlength{\parskip}{5pt}
\setlength{\emergencystretch}{3em}
\newcommand{\N}{\mathbb{N}}
\newcommand{\Z}{\mathbb{Z}}
\newcommand{\Q}{\mathbb{Q}}
\newcommand{\R}{\mathbb{R}}
\newcommand{\C}{\mathbb{C}}
\newcommand{\F}{\mathbb{F}}
\newcommand{\A}{\mathbb{A}}
\newcommand{\PP}{\mathbb{P}}
\newcommand{\E}{\mathbb{E}}
\newcommand{\eps}{\varepsilon}
\newcommand{\dd}{\,\mathrm{d}}
\begin{document}
\section*{UnsolvedMath ID: 800670}
\textbf{Source problem code:} catalogue entry
\subsection{Definitions and mathematical statement}
Let \(X\) be a topological space, and let \(C_p(X)\) be the set of continuous functions \(X\to\R\), equipped with the topology of pointwise convergence, inherited from \(\R^X\). The supplied excerpt does not specify the intended class of spaces \(X\). In particular, whether complete regularity and the Hausdorff separation axiom are assumed must be checked before the statement is finalized.

A topological space \(Y\) has the \emph{Reznichenko property} if, whenever \(A\subseteq Y\) and \(y\in\overline{A}\setminus A\), there exist pairwise disjoint nonempty finite sets \(F_n\subseteq A\), indexed by \(n\in\N\), such that every neighborhood of \(y\) meets \(F_n\) for all but finitely many \(n\).

An \emph{\(\omega\)-cover} of \(X\) is a family \(\mathcal U\) of subsets of \(X\) such that \(X\notin\mathcal U\) and every finite subset of \(X\) is contained in some member of \(\mathcal U\). It is open if all its members are open. Write \(\Omega(X)\) for the open \(\omega\)-covers admitted by the source's cover convention. The supplied excerpt does not specify whether admissible covers must be countable; this convention remains unresolved.

Under the usual interpretation of \(\omega\)-groupability, an \(\omega\)-cover \(\mathcal V\) is \emph{\(\omega\)-groupable} if it has a partition
\[
\mathcal V=\bigsqcup_{n\in\N}\mathcal V_n
\]
into nonempty finite families such that
\[
\forall K\subseteq X\text{ finite}\quad
\exists n_0\in\N\quad
\forall n\ge n_0\quad
\exists V\in\mathcal V_n\quad K\subseteq V.
\]
Let \(\Omega^{\mathrm{gp}}(X)\) denote the admissible open \(\omega\)-covers with this property. The selection statement
\[
\binom{\Omega}{\Omega^{\mathrm{gp}}}
\]
means explicitly that
\[
\forall\mathcal U\in\Omega(X)\quad
\exists\mathcal V\subseteq\mathcal U\quad
\mathcal V\in\Omega^{\mathrm{gp}}(X).
\]
An open \(\omega\)-cover \(\mathcal U\) is \emph{\(\omega\)-shrinkable} if there are closed subsets \(C_U\subseteq U\), one for each \(U\in\mathcal U\), such that \(\{C_U:U\in\mathcal U\}\) is an \(\omega\)-cover of \(X\); these closed sets need not be open.

\emph{Question, subject to the unresolved ambient and cover conventions.} For every space \(X\) in the source's intended class, is it true that
\[
C_p(X)\text{ has the Reznichenko property}
\quad\Longleftrightarrow\quad
\binom{\Omega(X)}{\Omega^{\mathrm{gp}}(X)}?
\]
The final displayed question in the supplied record has additional typographic decoration around \(\mathrm{gp}\). The display above uses the ordinary \(\Omega^{\mathrm{gp}}\) appearing in the preceding discussion; the record does not define a distinct property associated with that decoration.

\emph{Supplied context.} The record attributes to \([\mathrm{SakaiRez}]\) the equivalence between the Reznichenko property of \(C_p(X)\) and the assertion that every \(\omega\)-shrinkable open \(\omega\)-cover of \(X\) is \(\omega\)-groupable. It consequently presents the implication from \(\binom{\Omega}{\Omega^{\mathrm{gp}}}\) to the Reznichenko property as known, and asks about the converse.

The record also attributes to \([\mathrm{coc7}]\) the equivalence between the conjunction of countable fan tightness and the Reznichenko property for \(C_p(X)\), and the Hurewicz covering property for every finite power \(X^m\), \(m\ge1\). Here countable fan tightness means that, for every \(y\in Y\) and every sequence \((A_n)_{n\in\N}\) of subsets of \(Y\) with \(y\in\overline{A_n}\) for each \(n\), there are finite \(B_n\subseteq A_n\) such that \(y\in\overline{\bigcup_n B_n}\). The Hurewicz covering property can be expressed by requiring that, for every sequence \((\mathcal W_n)_{n\in\N}\) of open covers, finite families \(\mathcal H_n\subseteq\mathcal W_n\) can be chosen so that each point belongs to \(\bigcup\mathcal H_n\) for all but finitely many \(n\). These contextual equivalences are supplied source assertions whose ambient hypotheses have not been independently verified.
\subsection{Short English statement}
The supplied question asks whether the Reznichenko property of \(C_p(X)\) forces every admissible open \(\omega\)-cover of \(X\) to contain an \(\omega\)-groupable subcover. The reverse implication is presented as known in the source. The intended class of spaces, the countability convention for covers, and the typography of the final superscript require review.
\subsection{Sources}
\begin{itemize}
\item Local UnsolvedMath registry: supplied complete record with permanent ID \(800670\), exact source code \texttt{catalogue entry}, in the AMR Open Problem Lists collection. The registry lists Tsaban as proposer, gives the year \(2003\), and labels the problem open; this is a supplied status label, not an independently verified present-day assessment.
\item Tsaban, \emph{Selection Principles in Mathematics: A Milestone of Open Problems} (2003), Problem 43. Supplied source URL: \url{https://arxiv.org/abs/math/0312182}. The record reports access on \texttt{2026-10-16} and extraction by \texttt{arxiv\_latex}; these are supplied provenance statements, not a claim of access during this preparation.
\item Supplied attribution and context: \([\mathrm{SakaiRez},\ \text{Question }3.5]\) and \([\mathrm{reznicb}]\). The shrinkable-cover characterization is attributed to Sakai and \([\mathrm{SakaiRez}]\). Full bibliographic entries for these reference labels are absent from the supplied record.
\item The contextual equivalence involving countable fan tightness, the Reznichenko property, and the Hurewicz property of finite powers is attributed in the supplied excerpt to \([\mathrm{coc7}]\). Its full bibliographic entry and precise hypotheses are not supplied.
\end{itemize}
\end{document}
